% ================================================================
% ch05_prem_gaan.tex — अध्याय ५: प्रेम गान आ निरगुनी पद संग्रह
% ================================================================

\cleardoublepage
\section*{अध्याय ५: प्रेम गान आ निरगुनी पद संग्रह}
\addcontentsline{toc}{section}{अध्याय ५: प्रेम गान आ निरगुनी पद संग्रह}

\subsection*{५.१ प्रेम के स्वरूप आ विरह के वेदना}
\addcontentsline{toc}{subsection}{५.१ प्रेम के स्वरूप आ विरह के वेदना}

\textbf{मूल भोजपुरी रचना:}

\begin{quote}
\addfontfeatures{Color=3D3428}
पहले प्रेम पाछे सब कीजे, बिना प्रेम किस काम को जीजे।\\
प्रेम बिना जीवन है सूना, जैसे बिन जल मछली ऊना॥
\end{quote}

\textbf{भोजपुरी भावार्थ:} पहिले मन में सच्चा प्रेम पैदा करा --- प्रभु के प्रति, गुरु के प्रति, आ सत्य के प्रति। बाकी सब साधना आ उपाय बाद में करा। बिना ओह प्रेम के कवनो जीवन सार्थक नइखे। ऊ जीवन पानी बिना तड़पत मछली जइसन ह।

\begin{sloppypar}
{\engfont \textit{English Translation:} First cultivate love in your heart --- everything else follows.
Without that love, what is the worth of life? Life without love is empty --- like a fish writhing without water.}
\end{sloppypar}

\begin{center}
{\color{bordercolor}\rule{0.4\textwidth}{0.4pt}}
\end{center}

\subsection*{५.२ प्रेम के अग्नि आ भीतरी क्रांति}
\addcontentsline{toc}{subsection}{५.२ प्रेम के अग्नि आ भीतरी क्रांति}

\textbf{मूल भोजपुरी रचना:}

\begin{quote}
\addfontfeatures{Color=3D3428}
प्रेम अगन जब मन में जारे, तन की खाज सब भस्म निकारे।\\
काया ममता जरि जब जावे, तब दरिया रस पीवत धावे॥
\end{quote}

\textbf{भोजपुरी भावार्थ:} जब प्रेम के आग मन में जल उठेला त तन के सब मोह-ममता जला के भस्म हो जाला। जब शरीर के मोह जर जाला, तब आत्मा आनंद-रस पिए खातिर दौड़े लागेला।

\begin{sloppypar}
{\engfont \textit{English Translation:} When the fire of love ignites within the mind, all bodily attachments are burned to ash.
When ego is consumed, says Dariya, the soul races towards divine nectar.}
\end{sloppypar}

\begin{center}
{\color{bordercolor}\rule{0.4\textwidth}{0.4pt}}
\end{center}

\subsection*{५.३ सतगुरु-प्रेम आ पारख}
\addcontentsline{toc}{subsection}{५.३ सतगुरु-प्रेम आ पारख}

\textbf{मूल भोजपुरी रचना:}

\begin{quote}
\addfontfeatures{Color=3D3428}
प्रेम ज्ञान जब उपजे, चले जगत कंह झारी।\\
कहे दरिया सतगुरु मिले, पारख करे सुधारी॥
\end{quote}

\textbf{भोजपुरी भावार्थ:} जब मन में प्रेम आ ज्ञान दूनों एक साथ उपजेला त जीव संसार के भ्रम से दूर हो जाला। जब सतगुरु मिललें, उनकर नज़र से आत्मा के सुधार आ पारख भइल।

\begin{sloppypar}
{\engfont \textit{English Translation:} When love and knowledge arise together within, the soul transcends worldly illusion.
Says Dariya --- when the True Master is found, He tests and refines the soul with divine discernment.}
\end{sloppypar}

\begin{center}
{\color{bordercolor}\rule{0.4\textwidth}{0.4pt}}
\end{center}

\subsection*{५.४ गुरु-समर्पण आ परम पद}
\addcontentsline{toc}{subsection}{५.४ गुरु-समर्पण आ परम पद}

\textbf{मूल भोजपुरी रचना:}

\begin{quote}
\addfontfeatures{Color=3D3428}
मन दे गुरु को, तन दे गुरु को, जीव जान सब दे गुरु को।\\
कहे दरिया यह परखी बात, गुरु बिन मिले न परमात॥
\end{quote}

\textbf{भोजपुरी भावार्थ:} गुरु को मन, तन, जीव सब समर्पित कर दो। दरिया साहब कहत बाड़ें कि ई उनकर अनुभव कइल सत्य ह --- सतगुरु के बिना परमात्मा के दर्शन असंभव बा।

\begin{sloppypar}
{\engfont \textit{English Translation:} Offer your mind, body, life and soul to the True Master.
Says Dariya, this is a tested truth --- without the True Master, meeting the Supreme is impossible.}
\end{sloppypar}

\begin{center}
{\color{bordercolor}\rule{0.4\textwidth}{0.4pt}}
\end{center}

\subsection*{५.५ माया ठगिनी आ आडंबर खंडन}
\addcontentsline{toc}{subsection}{५.५ माया ठगिनी आ आडंबर खंडन}

\textbf{मूल भोजपुरी रचना:}

\begin{quote}
\addfontfeatures{Color=3D3428}
माया ठगिनी महा सयानी, जग को ठगत फिरत मनमानी।\\
जिन जान्यो तिन दूर निसारा, दरिया सत्य विचार हमारा॥
\end{quote}

\textbf{भोजपुरी भावार्थ:} माया बड़ी चालाक ठगनी ह, ऊ एह जगत के अपने मनमाने तरीके से ठगत रहेली। जे एकरा के जान लेला उ एकरा से दूर हो जाला — दरिया साहब कहत बाड़ें कि ई हमर सच्चा विचार ह।

\begin{sloppypar}
{\engfont \textit{English Translation:} Maya (illusion) is a supremely cunning trickster, deceiving the world in her own whimsical ways.
One who truly knows her nature breaks free from her grasp --- this is Dariya's essential truth.}
\end{sloppypar}

\begin{center}
{\color{bordercolor}\rule{0.4\textwidth}{0.4pt}}
\end{center}

\subsection*{५.६ नाम-सुमिरन आ अनहद नाद}
\addcontentsline{toc}{subsection}{५.६ नाम-सुमिरन आ अनहद नाद}

\textbf{मूल भोजपुरी रचना:}

\begin{quote}
\addfontfeatures{Color=3D3428}
सुमिरन ऐसा कीजिए, दूजा मन नहिं आय।\\
रोम रोम में नाम रमै, सुरति निरति ठहराय॥
\end{quote}

\textbf{भोजपुरी भावार्थ:} नाम के सुमिरन अइसन गहरा होखे कि मन में कवनो दूसर विचार न आवे। नाम रोम-रोम में बस जाव आ ध्यान पूरी तरह स्थिर हो जाव।

\begin{sloppypar}
{\engfont \textit{English Translation:} Remember the Name so deeply that no second thought enters the mind.
May the Name pervade every pore of your being, stilling awareness into perfect stillness.}
\end{sloppypar}

\begin{center}
{\color{bordercolor}\rule{0.4\textwidth}{0.4pt}}
\end{center}

\subsection*{५.७ साध के लक्षण आ नम्रता}
\addcontentsline{toc}{subsection}{५.७ साध के लक्षण आ नम्रता}

\textbf{मूल भोजपुरी रचना:}

\begin{quote}
\addfontfeatures{Color=3D3428}
दरिया लच्छन साध का, क्या गिरही क्या भेख।\\
नि:कपटी निरसंक रहि, बाहर भीतर एक॥
\end{quote}

\textbf{भोजपुरी भावार्थ:} सच्चे साधु के लक्षण ह कि ऊ निष्कपट, निर्भय आ बाहर-भीतर एक होखे। कवनो दोहरापन न होखे।

\begin{sloppypar}
{\engfont \textit{English Translation:} The true saint is free from deceit, fearless, and the same within as without --- no duality, no pretense.}
\end{sloppypar}

\begin{center}
{\color{bordercolor}\rule{0.4\textwidth}{0.4pt}}
\end{center}

\subsection*{५.८ सहज समर्पण आ मोक्ष}
\addcontentsline{toc}{subsection}{५.८ सहज समर्पण आ मोक्ष}

\textbf{मूल भोजपुरी रचना:}

\begin{quote}
\addfontfeatures{Color=3D3428}
दरिया सरल समर्पणा, यही मोक्ष के उपाय।\\
जो कोई इस राह चले, सो जग से उठि जाय॥
\end{quote}

\textbf{भोजपुरी भावार्थ:} दरिया साहब कहत बाड़ें कि निष्काम भाव से कइल गइल सहज आ सरल समर्पण ही मोक्ष के एकमात्र रास्ता ह। जे एह पर चलेला उ संसार चक्र से मुक्त हो जाला।

\begin{sloppypar}
{\engfont \textit{English Translation:} Says Dariya, simple and sincere surrender is the only path to liberation.
Whoever walks this road rises above the world and frees the soul from rebirth.}
\end{sloppypar}
